% EM structure and teaching revision 20260928
\subsection{复功率守恒与无功补偿}
对线性稳态网络，逐支路的复功率相加可同时检验有功与无功的端口平衡。以下沿用前文的 $Q_r$ 表示无功功率，避免与电荷量混淆。
在感性负载两端并接电容，保持电压有效值不变时，电容提供 $Q_{r,C}=-\omega CU^2$。若负载从功率因数角 $\varphi_1$ 改善到 $\varphi_2$，所需
\[
 C=\frac{P(\tan\varphi_1-\tan\varphi_2)}{\omega U^2}.
\]
补偿减少电源与线路中的无功电流，不会把负载原来需要的有功功率消除。过度补偿会使输入变成容性，也不能用“电流越小越好”取代具体设计条件。
理想变压器保持有功和无功的复功率关系，负载折算必须乘匝数比的平方；电压按匝数比变化，而电流反比变化。若结果让无损变压器无端产生实功率，应检查电流参考方向或折算比是否倒置。

